% Options for packages loaded elsewhere
\PassOptionsToPackage{unicode}{hyperref}
\PassOptionsToPackage{hyphens}{url}
\documentclass[]{article}
\usepackage{amsmath,amssymb}
\usepackage{iftex}
\ifPDFTeX
  \usepackage[T1]{fontenc}
  \usepackage[utf8]{inputenc}
  \usepackage{textcomp}
\else
  \usepackage{unicode-math}
  \defaultfontfeatures{Scale=MatchLowercase}
  \defaultfontfeatures[\rmfamily]{Ligatures=TeX,Scale=1}
\fi
\usepackage{lmodern}
\IfFileExists{upquote.sty}{\usepackage{upquote}}{}
\IfFileExists{microtype.sty}{\usepackage[]{microtype}\UseMicrotypeSet[protrusion]{basicmath}}{}
\makeatletter
\@ifundefined{KOMAClassName}{\IfFileExists{parskip.sty}{\usepackage{parskip}}{\setlength{\parindent}{0pt}\setlength{\parskip}{6pt plus 2pt minus 1pt}}}{\KOMAoptions{parskip=half}}
\makeatother
\usepackage{xcolor}
\setlength{\emergencystretch}{3em}
\providecommand{\tightlist}{\setlength{\itemsep}{0pt}\setlength{\parskip}{0pt}}
\setcounter{secnumdepth}{-\maxdimen}
\ifLuaTeX\usepackage{selnolig}\fi
\usepackage{bookmark}
\IfFileExists{xurl.sty}{\usepackage{xurl}}{}
\urlstyle{same}
\hypersetup{hidelinks,pdfcreator={LaTeX via pandoc}}
\author{}
\date{}

\begin{document}

\section{Relational Emergence Model (REM)}
\subsection{Core Specification v1.1}

\section{0. Core Statement}
REM treats tensor-product structure of Hilbert space as a variable structure selected via competing informational and dynamical constraints.

REM is a structural selection framework, not a modification of quantum dynamics.

\section{1. Three-Layer Architecture}
\subsection{Layer 0 --- Pre-relational description}
Hilbert space H is considered without specifying a preferred tensor-product decomposition.

Layer 0 does not represent a physical substrate. It denotes absence of articulated relational structure.

\subsection{Layer 1 --- Relational articulation}
Differentiation introduces an explicit tensor structure:

D : H → H\_S ⊗ H\_O

A quantum state admits Schmidt representation:

\textbar Ψ⟩ → Σ c\_i \textbar s\_i⟩ ⊗ \textbar o\_i⟩

Relational structure becomes explicit.

\subsection{Layer 2 --- Phenomenal stabilization}
Environmental decoherence stabilizes articulated correlations.

Ordering relation:

D → E

\section{2. Differentiation map}
D is a structural selection map.

D is not a dynamical law and is not a unitary operator.

D is defined implicitly through a selection problem over candidate factorizations.

Variational principles provide tractable models for D but do not exhaust its possible formalizations.

\section{3. Factorization space}
Candidate tensor decompositions:

F = \{F\_k\}

Finite-dimensional approximation:

F ≅ U(N)/(U(n\_A) ⊗ U(n\_B))

This representation is heuristic and applies primarily in finite-dimensional Hilbert spaces.

\section{4. Variational ansatz}
Lowest-order functional:

Φ(F) = Φ\_S(F) + λ Φ\_H(F)

Higher-order generalization:

Φ(F) = Φ\_S + λ Φ\_H + Σ μ\_i Φ\_i

Optimal factorization:

F* = argmax\_F Φ(F)

The functional provides an effective model for relational selection.

\section{5. Informational functional}
Φ\_S = I(A:B)

I(A:B) = S(ρ\_A) + S(ρ\_B) − S(ρ\_AB)

Von Neumann entropy:

S(ρ) = −Tr(ρ log ρ)

Units may be nats or bits depending on logarithm base.

\section{6. Dynamical functional}
Φ\_H measures dynamical coupling across candidate partitions.

Lowest-order approximation:

Φ\_H = −Tr(ρ H\_boundary)

H\_boundary represents interaction terms spanning the factorization boundary.

Interpretation: dynamical self-containment measure.

\section{7. Interpretation of λ}
Primary meaning: relative weighting between informational articulation and dynamical self-containment.

Secondary interpretations include a thermodynamic analogy, an emergent scale parameter, and a timescale relation λ \textasciitilde{} τ\_D / τ\_dyn.

λ is an effective parameter, not a universal constant.

\section{8. Crossover structure}
λ* is defined by equality of competing factorization scores.

λ* is model dependent.

REM predicts existence of crossover regimes, not universal λ* values.

\section{9. Role of toy models}
Small systems illustrate nontrivial structure in factorization landscape.

Toy models serve as existence proofs. They do not claim empirical completeness.

\section{10. Relation to prior work}
REM is situated within discussions concerning the non-uniqueness of subsystem structure in quantum theory.

Zanardi, Lidar, and Whaley demonstrated that tensor product structure is not uniquely determined by Hilbert space alone. Cotler, Penington, and Ranard showed that preferred factorizations may arise through optimization. Carroll and Singh proposed a dynamical quantum-mereological framework. Donnelly and Freidel showed that subsystem structure in gauge theories is nontrivial. Rovelli introduced relational quantum mechanics, while Zurek showed that decoherence selects stable pointer observables once subsystem structure is assumed.

REM contributes a complementary structural proposal: tensor-product structure itself is treated as a variable object selected through a variational competition between informational articulation and dynamical self-containment.

\section{11. Scope of theory}
REM addresses the structural selection problem.

REM does not introduce new dynamical laws, hidden variables, or modifications to quantum mechanics.

REM provides structural constraints on subsystem articulation.

\section{12. Terminology}
Differentiation: selection of relational tensor structure.

Factorization: tensor product decomposition.

Articulation: explicit correlation structure.

Self-containment: weak dynamical coupling across partition.

Crossover: change in preferred factorization regime.

\section{13. Extension directions}
Multi-qubit scaling, tensor-network articulation, continuous Hilbert spaces, renormalization structure, and quantum-channel formulations of differentiation.

\section{14. Status of theory}
REM is a structured foundational proposal.

It provides a mathematically formulated research program for relational structure selection.

Future work will determine empirical and mathematical generality.

\end{document}
